\centerline{\bf \Huge Входная контрольная работа}
\centerline{\bf \Huge }

\begin{flushright}
        \scriptsize{\textbf{20:31 прибыл Годжу Сатору}}
\end{flushright}

\problem{1-\textbf{3!}} Сформулируйте признаки делимости на:

а-\textbf{1!}) 5

б-\textbf{1!}) 3 и 9

в-\textbf{1!}) 4 и 8

\problem{2-\textbf{6!}} Сформулируйте \textbf{ВСЕ} признаки и свойства параллельных прямых и секущей.

\problem{3-\textbf{3}} Чему равны следующие выражения:

а-\textbf{1}) $2\cdot1,57\cdot0,43 + 1,57^2 + 0,43\cdot0,43=?$

б-\textbf{2}) $3\cdot3,99\cdot3,49^2 + 3,99^2\cdot3,99 - 3\cdot3,99^2\cdot3,49 - \dfrac{3,49^5}{3,49^2}$

\problem{4-\textbf{3}} В треугольниках $ABC$ и $A_1B_1C_1$: $AB=A_1B_1$, $\angle A=\angle A_1$, $\angle B=\angle B_1$, точки $D$ и $D_1$ лежат соответственно на сторонах $AC$ и $A_1C_1$, причем $CD=C_1D_1$. Докажите, что $\triangle BDC = \triangle B_1D_1C_1$. Сравните отрезки $BD$ и $B_1D_1$.

\problem{5-\textbf{6}} Прямые $y=x+3$ \ и \ $y=2x-1$ пересекаются в точке $A$, а прямые $y=7-3x$ \ и \ $y=5$ в точке $B$. Найдите уравнение прямой, проходящей через точки $A$ и $B$.

\problem{6-\textbf{6}} В треугольниках $ABC$ и $A_1B_1C_1$: $BC=B_1C_1$, $\angle C=\angle C_1$ и $AB + AC = A_1B_1 + A_1C_1$, $BD$ и $B_1D_1$ - медианы этих треугольников. Докажите, что $BD=B_1D_1$.



\begin{flushright}
        \scriptsize{\textbf{Доржеев Анджур Аркадьевич}}
\end{flushright}